
@@M0@@
où @@M1@@ est un idéal homogène de @@M2@@, de
degré @@M3@@. Il résulte alors de la définition (4.9.1) que @@M4@@ est le séparé complété de @@M5@@ pour la topologie définie par les @@M6@@. Cette topologie est plus fine que la topologie définie par le poids; nous allons montrer qu'elle lui est équivalente, ce qui entrainera la proposition.

Posons @@M7@@, et soit @@M8@@ la @@M9@@-algèbre engendrée par des générateurs @@M10@@ soumis aux relations :
@@M11@@
où les @@M12@@ désignent les fonctions symétriques élémentaires. D'après 4.4, l'homomorphisme canonique @@M13@@ est injectif. On obtient :
@@M14@@
Comme @@M15@@ est un polynome, il en est de même pour @@M16@@ pour tout @@M17@@. Par suite, les @@M18@@ sont nilpotents, et il existe un entier @@M19@@ tel que tout polynome de degré @@M20@@ par rapport aux @@M21@@, à coefficients dans @@M22@@, soit nul. Par conséquent, tout polynome de @@M23@@, de poids @@M24@@, est nul, i.e. tout polynome de @@M25@@ de poids @@M26@@ appartient à @@M27@@. Il s'ensuit que les deux topologies considérées sont équivalentes, ce qui achève la démonstration.


\section*{5. Applications aux fibrés projectifs; étude de \texorpdfstring{@@M28@@}{f*}}

5.1.
Dans ce paragraphe, @@M29@@ est un schéma quasi-compact, @@M30@@ un @@M31@@-Module localement libre de rang @@M32@@, @@M33@@ le fibré projectif associé à @@M34@@. Sur @@M35@@, on a la suite exacte canonique :
@@M36@@
